% Analyse qualité du dépôt ThesisManager (mod_these, mod_thesesynthesis), suivi de l'analyse du
% 25 octobre 2026.
% Compilation : pdflatex analyse-qualite-2026-10-30.tex (deux passes), paquets TeX Live
% latex-recommended, latex-extra, lang-french.
\documentclass[11pt,a4paper]{article}
\usepackage[T1]{fontenc}
\usepackage[utf8]{inputenc}
\usepackage{lmodern}
\usepackage[french]{babel}
\usepackage[margin=2.2cm]{geometry}
\usepackage{booktabs}
\usepackage{longtable}
\usepackage{tabularx}
\usepackage{array}
\usepackage{enumitem}
\usepackage{xcolor}
\usepackage[colorlinks=true,linkcolor=blue!50!black,urlcolor=blue!50!black]{hyperref}

\setlist{nosep}
\sloppy
\setlength{\emergencystretch}{3em}
\newcommand{\fichier}[1]{\texttt{\small #1}}
\newcommand{\code}[1]{\texttt{#1}}
\newcommand{\sev}[1]{\textbf{#1}}
\newcommand{\corrige}{\textcolor{green!45!black}{\textbf{corrigé}}}
\newcommand{\partiel}{\textcolor{orange!80!black}{\textbf{partiel}}}
\newcommand{\present}{\textcolor{red!70!black}{\textbf{présent}}}
\newcolumntype{Y}{>{\raggedright\arraybackslash}X}

\title{ThesisManager : analyse qualité du projet\\
  \large \code{mod\_these} 0.23.1 (2026103000) et \code{mod\_thesesynthesis} 0.7.2 (2026102901)}
\author{Sébastien Lefebvre}
\date{30 octobre 2026 --- commit \code{5ccadce} (fusion de la demande n\textsuperscript{o}~31)}

\begin{document}
\maketitle

\section*{Synthèse}

En cinq jours, le dépôt est passé de 0.13.0 à 0.23.1 : 79 commits, 21 demandes de fusion, +30~\%
de code (45~900 lignes de PHP), dix classes de domaine nouvelles, deux tables, onze étapes de
mise à niveau. La tenue générale reste remarquable : \code{phpcs} sans erreur ni avertissement,
aucune erreur de syntaxe, parité parfaite des 1~112 chaînes français/anglais, 233 tests PHPUnit
(contre 183) et 16 scénarios Behat (contre 9), \fichier{db/upgrade.php} cohérent avec
\fichier{install.xml}, \code{main} verte. Toutes les pages nouvelles contrôlent connexion, droits,
appartenance du dossier à l'activité et \code{sesskey}.

Le constat de sécurité élevé de l'analyse précédente (rapport et grille lisibles par l'étudiant)
est \corrige, ainsi que le calendrier des soutenances, les noms de personnes dans les valeurs par
défaut et le texte d'aide de l'embargo. En revanche, la quasi-totalité des autres constats est
encore présente : la vitesse de développement fonctionnel a primé sur la résorption de la dette.
Par ordre d'importance :

\begin{enumerate}
  \item \sev{Sécurité (moyen, aggravé)} : la demande de signature du directeur extérieur part
    toujours, avant toute validation, vers l'adresse saisie par l'étudiant ; la nouvelle action
    groupée « relancer » la renvoie avec un jeton neuf, et « Mes signatures » permet de signer et
    de prévisualiser sans jeton sur simple égalité d'adresse (§\ref{sec:securite}).
  \item \sev{Droits} : \code{managedeve} ne porte toujours pas \code{RISK\_XSS} alors qu'elle
    permet d'éditer et désormais d'importer le HTML des modèles ; la substitution des balises de
    courriel reste faite en plusieurs passes.
  \item \sev{Processus} : deux fusions sur huit ont encore rendu \code{main} rouge (demandes
    n\textsuperscript{os}~25 et 28) ; la protection de branche n'est toujours pas active
    (§\ref{sec:processus}).
  \item \sev{Données} : sauvegarde (\code{templatehash}, \code{validatedby} non annoté,
    remappage entre sites), métadonnées et requêtes RGPD incomplètes, politique de conservation du
    procès-verbal absente, index manquants (§\ref{sec:donnees}).
  \item \sev{Architecture} : la dette mesurée a augmenté : 715 appels aux façades depuis
    \fichier{classes/} (606), 20 méthodes de plus de 80 lignes (12), requêtes par ligne dans le
    suivi multi-promotions (§\ref{sec:architecture}).
\end{enumerate}

\section{Méthode et périmètre}

L'analyse porte sur l'arbre de travail au commit \code{5ccadce} (identique à \code{origin/main}),
hors \fichier{mod/these/thirdparty}. Elle combine :
\begin{itemize}
  \item l'exécution locale de \code{php -l} (PHP 8.3) sur tous les fichiers et de
    \code{phpcs --standard=phpcs.xml} avec \code{moodlehq/moodle-cs} 3.7 : aucun écart ;
  \item l'état des huit dernières exécutions de la CI sur \code{main} (n\textsuperscript{os}~71
    à 93) ;
  \item un script fondé sur le \emph{tokenizer} PHP, appliqué à l'arbre actuel et à l'arbre
    0.13.0 (\code{a0e555c}), pour des mesures comparables d'une analyse à l'autre ;
  \item trois lectures ciblées : sécurité, architecture, puis tests, données, CI et
    documentation ; chaque constat de l'analyse du 25 octobre est requalifié (\corrige,
    \partiel, \present) et les constats marquants ont été revérifiés dans le source.
\end{itemize}
Aucun test PHPUnit ni Behat n'a été exécuté ici (pas de Moodle installé) ; la CI les exécute et
ils passent sur \code{main}.

\section{Chiffres clés}

\begin{center}
\begin{tabular}{lrr}
\toprule
 & \textbf{0.13.0 (25/10)} & \textbf{0.23.1 (30/10)} \\
\midrule
PHP hors bibliothèque tierce & 35~150 lignes & 45~860 lignes \\
dont tests & --- & 10~400 lignes \\
Chaînes de langue \code{en} / \code{fr} (les deux modules) & 912, 0 écart & 1~159, 0 écart \\
Façades \code{these\_*()} (\fichier{locallib/}) & 240 & 244 \\
Appels aux façades / appels directs aux classes & 1~610 / 34 & 1~964 / 557 \\
Appels aux façades depuis \fichier{classes/} & 606 & 715 \\
Classes \fichier{classes/local/} (100~\% statiques) & 23 (264 méthodes) & 33 (405 méthodes) \\
Méthodes de plus de 80 lignes & 12 & 20 \\
Tests PHPUnit & 183 & 233 (209 + 24) \\
Scénarios Behat (dont JavaScript) & 9 (1) & 16 (2) \\
Gabarits Mustache, tous référencés & 33 + 8 & 37 + 9 \\
Tables & 12 + 2 & 14 + 2 \\
Écarts \code{phpcs} & 30 & 0 \\
Images du guide (non citées) & 78 (11), 13~Mo & 97 (11), 15~Mo \\
\bottomrule
\end{tabular}
\end{center}

\section{État du dépôt et processus de livraison}\label{sec:processus}

\paragraph{CI.}
\code{main} est verte depuis la demande n\textsuperscript{o}~29. Mais les fusions
n\textsuperscript{os}~25 et 28 ont produit une exécution en échec sur \code{main} : la règle
« CI verte et relecture avant fusion » de CONTRIBUTING n'est toujours pas imposée par une
protection de branche. C'est la mesure la moins coûteuse et la plus efficace du plan d'action.

\paragraph{Publication.} \partiel. La première section de \fichier{CHANGES.md} porte bien la
version de \code{version.php} (0.23.1), mais rien ne le vérifie : le job « Publication »
extrait toujours la première section sans contrôle. \code{mod\_thesesynthesis} n'a pas
d'étiquette propre ; ses versions ne se lisent que dans les titres de \fichier{CHANGES.md}.

\paragraph{Gouvernance.} \present. Pas de \fichier{LICENSE} (exigé par le répertoire des
plugins Moodle), ni de \fichier{SECURITY.md}, ni de \fichier{CODEOWNERS}.

\paragraph{Rythme.} 79 commits en cinq jours, jusqu'à sept montées de version mineure. Les
commits restent bien découpés, mais ce rythme laisse peu de place à la relecture humaine que
CONTRIBUTING exige, et la maturité \code{ALPHA} de \code{mod\_these} doit être maintenue tant
que les constats moyens de sécurité ne sont pas traités.

\section{Sécurité}\label{sec:securite}

\subsection*{Ce qui est fait correctement}

Tout ce que l'analyse précédente relevait est maintenu, et le code nouveau suit les mêmes règles :
\code{require\_login} et contrôle de droit sur les quatorze pages ajoutées ou refondues
(\fichier{daterequest.php}, \fichier{newdate.php}, \fichier{corrections.php}, \fichier{a6.php},
\fichier{a6import.php}, \fichier{bulk.php}, \fichier{griditems.php},
\fichier{settingstransfer.php}, \fichier{import.php}, \fichier{scan.php}, \fichier{refuse.php},
\fichier{report.php}, \fichier{reporttemplate.php}, \fichier{questionnaire.php}) ; dossier chargé
avec \code{theseid} ; actions modifiantes par \code{moodleform} ou \code{confirm\_sesskey} ; SQL
toujours paramétré ; aucun \code{unserialize}, \code{eval}, \code{exec} ni accès direct à
\code{\$\_GET} ; triple moustache alimentée seulement par \code{format\_text} ou des rendus ;
téléchargement forcé des fichiers. Le transfert de paramètres valide modèles et images et borne
les activités sources à celles où l'utilisateur est DEVE. \fichier{sign.php} exige désormais le
jeton.

\subsection*{Suivi des constats du 25 octobre}

\begin{longtable}{@{}p{4.6cm}p{1.6cm}p{9.3cm}@{}}
\toprule
\textbf{Constat} & \textbf{État} & \textbf{Preuve} \\
\midrule
\endhead
Rapport et grille dans la zone \code{generated} (élevé) & \corrige &
  \fichier{classes/local/documents.php} 126--133 ne copie que les documents non
  \code{staffonly} ; la grille n'est rendue que pour l'archive de l'université ; le rapport reste
  dans \code{directorreport}, filtré par \fichier{rights.php} 166--168. \\
Signature du directeur extérieur contrôlée par l'étudiant & \present{} (aggravé) &
  Demande émise à la soumission (\fichier{dossiers.php} 431--433,
  \fichier{subjectchange.php} 163--165) ; signature sans jeton par égalité d'adresse
  (\fichier{signatureagent.php} 367--368, \fichier{mysignatures.php} 122--139) ; renvoi par
  l'action groupée (\fichier{classes/local/bulk.php} 72--81). \\
Identité par égalité d'adresse & \present &
  \fichier{rights.php} 239--252, \fichier{jury.php} 787--799, \fichier{signatureagent.php} 368,
  \fichier{thesesynthesis/locallib.php} 330--338 ; aucun \code{userid} stocké, aucun contrôle
  \code{suspended}/\code{deleted}. \\
Sujet non masqué dans certains courriels & \corrige{} (requalifié) &
  La plupart des envois passent le dossier. Les envois encore faits sans dossier n'exposent le
  titre qu'à des destinataires qui doivent le connaître : convocations et archive de
  l'université (\fichier{documents.php} 226 et 260), désignation du jury (\fichier{jury.php}
  742), demande de signature aux signataires de l'université et au directeur extérieur
  (\fichier{signatures.php} 1070, \fichier{sign.php} 139). Il n'y a pas de fuite ; passer tout
  de même le dossier à ces appels, par cohérence, pour qu'un futur destinataire ne reçoive pas le
  titre par défaut. \\
Double substitution des balises & \present &
  \code{str\_replace} en tableau toujours utilisé (\fichier{notifications.php} 184--191) : une
  valeur substituée tôt contenant \code{\{\{...\}\}} est elle-même substituée. \\
\code{RISK\_XSS} sur \code{managedeve} ; balises \code{<tcpdf>} & \present &
  \fichier{db/access.php} 82 ; \fichier{classes/pdf/html\_document.php} 134--142 ne retire que les
  images. Le transfert de paramètres permet désormais d'\emph{importer} ce HTML. \\
Constats faibles & \present &
  Noms non échappés dans les redirections (\fichier{jurynotify.php} 47, \fichier{review.php} 209,
  \fichier{documents.php} 84) ; énumération par l'import CSV ; formules dans les exports
  (\fichier{export.php} 80, \fichier{planning\_page.php} 119) ; jeton dans
  \code{\$PAGE->set\_url} (\fichier{jury.php} 57, \fichier{sign.php} 55 et 113) ; annuaire
  global modifiable par toute DEVE. \\
\bottomrule
\end{longtable}

\subsection*{Nouveaux constats}

\begin{description}[style=nextline]
  \item[Moyen --- l'action groupée étend le constat sur le directeur extérieur.]
    « Relancer » (\fichier{classes/local/bulk.php} 72--81) émet un jeton neuf, valable 60 jours,
    pour chaque signataire extérieur de chaque dossier coché, adresse \code{extdirectoremail}
    saisie par l'étudiant comprise. Un étudiant qui a saisi une adresse qu'il contrôle reçoit un
    lien de signature à chaque relance de la DEVE.
  \item[Faible --- prévisualisation dans « Mes signatures ».]
    \fichier{mysignatures.php} 56--73 sert le PDF complet de toute signature dont l'adresse
    correspond au compte connecté, sans jeton : conséquence du constat sur l'identité par
    égalité d'adresse.
  \item[Faible --- import CSV vers des comptes non inscrits.]
    \fichier{classes/local/import.php} 265--300 rattache un dossier, voire un directeur, à
    n'importe quel compte du site (adresse ou \code{idnumber}), sans \code{is\_enrolled},
    contrairement à l'import A6 (\fichier{a6.php} 529).
  \item[Faible --- fichiers tableur sans limite.]
    \fichier{classes/local/a6.php} 437--468 charge entièrement le classeur (archive zip) par
    PhpSpreadsheet, sans \code{maxbytes} ni plafond de lignes : un fichier piégé peut épuiser la
    mémoire (réservé à la DEVE).
  \item[Faible --- dépôts sans restriction de type.]
    \fichier{classes/local/files.php} 41--45 (\code{accepted\_types '*'}) : le téléchargement
    forcé empêche le XSS direct, mais des fichiers \code{.html} ou \code{.svg} sont transmis au
    jury. Limiter le manuscrit au PDF.
  \item[Faible --- transfert de paramètres.]
    \code{section} des éléments de grille importée sans contrôle
    (\fichier{settingstransfer.php} 198).
\end{description}

\section{Architecture et maintenabilité}\label{sec:architecture}

\paragraph{Amélioré.}
Les appels directs aux classes passent de 34 à 557 (dont 333 par \code{use}) : les pages
commencent à contourner les façades. Plus aucune classe de \fichier{classes/local/} ne fait
\code{require\_once locallib.php}. \code{these\_deve\_agents} n'est plus mort. Aucun \code{TODO},
\code{print\_error} ni API dépréciée. Gabarits tous référencés, parité des langues parfaite.

\paragraph{Inchangé.}
Les 244 façades restent l'API principale (1~008 appels hors tests) et les classes de domaine
continuent de les consommer : la dépendance circulaire classe~$\rightarrow$~façade~$\rightarrow$
classe s'étend (715 appels depuis \fichier{classes/}, dont 228 vers les façades du même
domaine). Neuf constantes métier vivent toujours dans \fichier{locallib/}.
\code{mod\_thesesynthesis} n'a toujours pas de \fichier{classes/local/} (\fichier{locallib.php} :
528 lignes, 14 fonctions ; \code{thesesynthesis\_search}, 94 lignes, imbrication 6).
\code{build\_document\_context} reste à 372 lignes. Années d'étude dupliquées
(\fichier{subject\_form.php} 49, \fichier{doctemplates.php} 649). \code{catch} FPDI muet
(\fichier{doctemplates.php} 814).

\paragraph{Dégradé ou nouveau.}
\begin{itemize}
  \item Classes plus larges : \fichier{signatures.php} 1~258 lignes et 47 méthodes,
    \fichier{doctemplates.php} 1~143, \fichier{dossiers.php} 953, \fichier{jury.php} 850 (+64~\%).
    20 méthodes de plus de 80 lignes (12), dont trois contrôleurs \code{*\_page::output} ; six
    contrôleurs HTTP dans l'espace de noms de domaine (cinq).
  \item Code mort : \code{these\_render\_file\_links}, \code{these\_render\_actions},
    \code{these\_render\_progress}, \code{these\_sign\_from\_form} (façade et méthode),
    \code{these\_import\_columns} ; sept façades ne servent qu'aux tests ; deux chaînes inutilisées
    (\code{closedossier}, \code{ba6notyet}).
  \item Performance : \code{thesesynthesis\_get\_dossiers} fait désormais deux requêtes
    \code{user} et un appel à \code{these\_user\_knows\_subject} par ligne
    (\fichier{thesesynthesis/locallib.php} 170, 173, 189) ; même motif dans
    \code{thesesynthesis\_search} (473--474). Sur plusieurs promotions, la page du suivi coûte
    plusieurs centaines de requêtes.
  \item Dépendance inversée accrue : huit fichiers de \code{mod\_these} lisent tables,
    configuration ou capacités de \code{mod\_thesesynthesis} (trois auparavant), dont
    \fichier{scopesettings.php} 163 sans garde de version.
  \item Erreurs perdues : \fichier{backlinks.php} 105 renvoie \code{null} sans trace,
    \fichier{a6import.php} 62 remplace l'exception par un message générique sans journalisation.
  \item Amorce de page recopiée dans 38 pages sur 46 (29) ; prologue des formulaires dans 29
    classes (22).
  \item \fichier{db/upgrade.php} 39--45 appelle le code courant (\code{these\_document\_definitions},
    \code{these\_document\_filename}) : une refonte future cassera la montée depuis les anciennes
    versions. Figer les valeurs dans l'étape.
\end{itemize}

\section{Tests}

\paragraph{Progrès.}
233 tests (+50), environ 1~540 assertions, 16 scénarios Behat dont un dépôt de fichier ; test
du calendrier, de la nouvelle date, de la révocation des signatures au renvoi, du double passage
des relances.

\paragraph{Zones toujours minces.}
\begin{itemize}
  \item Sauvegarde/restauration : un seul test (\fichier{backup\_restore\_test.php} 79), sans
    signatures, scellés, grille ni questions ; rien pour \code{mod\_thesesynthesis}.
  \item Confidentialité : cinq tests, ni signature enregistrée, ni \code{these\_a6}, ni
    descellement (\fichier{provider.php} 546).
  \item Formulaires : 23 \code{validation()}, deux testées.
  \item Import CSV : ni encodage, ni BOM, ni séparateur, ni guillemets.
  \item Sécurité : aucun test pour les constats ouverts (relance du directeur extérieur,
    signature sans jeton).
  \item Environ 50 assertions sur du HTML exact ; \code{@covers} en docblock dans les 71 fichiers
    (attributs \code{\#[CoversClass]} attendus par PHPUnit 11) ; \code{upgrade\_test.php} : un
    seul test pour douze étapes ; aucun Behat pour \code{mod\_thesesynthesis}.
\end{itemize}

\section{Base de données, sauvegarde, confidentialité, tâches}\label{sec:donnees}

\paragraph{Schéma.}
\fichier{db/upgrade.php} (douze étapes, de 2026102504 à 2026102901) est cohérent avec
\fichier{install.xml} : tables et champs ajoutés présents, champs retirés absents, défaut de
\code{minjury} traité, savepoints dans l'ordre, aucune version datée dans le futur. Restent
\present : index sur \code{these\_log.userid}, \code{these\_signature.userid},
\code{these\_dossier.defensedate}, \code{these\_subjecthistory.directorid} et \code{changedby} ;
recherches par \code{LOWER(email)} ; clé étrangère \code{these\_question.theseid} ;
\code{these\_juror.email} à 100 caractères contre 255.

\paragraph{Sauvegarde et restauration.}
Les colonnes et tables nouvelles sont incluses (\code{these\_setting}). Restent \present :
\code{templatehash} absent (\fichier{backup\_these\_stepslib.php} 100--103) ; \code{jurorid} et
\code{questionid} conservés sans \code{is\_samesite} (\fichier{restore\_these\_stepslib.php} 243
et 164) ; dossiers ignorés sans trace (219--221) ; règles \code{STAGESYNTHESIS*}. Nouveau :
\code{validatedby} est remappé (284) mais pas annoté (149--157), il retombe à 0 sur un autre
site ; \code{these\_a6}, table de site comme l'annuaire, n'est sauvegardée nulle part, sans que ce
soit documenté.

\paragraph{API de confidentialité.} \partiel.
Métadonnées ajoutées pour \code{these\_juror}, \code{these\_a6} et \code{validatedby}. Toujours
absents : directeur extérieur (nom, grade, établissement, adresse), \code{correctionsnote},
\code{daterequestcomment}, \code{eligibleby}, \code{accesstoken}, \code{token},
\code{contenthash}, \code{verifiedby}, historique de sujet ; ces colonnes manquent aussi aux
requêtes de contexte, de liste et de suppression. Seul le jury est déclaré comme destination
externe. Surtout, la politique de conservation du procès-verbal n'est toujours pas écrite : la
suppression d'un étudiant efface le PV scellé, celle d'un membre du personnel le descelle.

\paragraph{Tâches.}
\code{send\_reminders} reste idempotente et testée. La nouvelle \code{send\_digest} (le lundi,
\fichier{db/tasks.php} 39--46) n'enregistre aucun horodatage d'envoi
(\fichier{notifications.php} 667--698) : une exécution manuelle ou une reprise renvoie le
récapitulatif. Journalisation conditionnelle dans les deux tâches. Calendrier : \corrige{}
(\fichier{classes/local/calendar.php}, testé). Toujours \code{email\_to\_user} sans
\fichier{db/messages.php}, et aucun événement Moodle.

\section{CI/CD et outillage}

Aucun des points de l'analyse précédente n'a été traité : \code{cancel-in-progress}
inconditionnel (y compris sur \code{main}) ; \code{contents: write} sur tout le job « Guide PDF »
et publication sur étiquette avec \fichier{CHANGES.md} entier ; \code{moodle-cs} non épinglé ;
pas d'étape \code{grunt}, ni \fichier{.map}, ni \fichier{Gruntfile.js} ; couverture dans le
journal seulement ; pas de \code{timeout-minutes} ; actions épinglées par étiquette majeure ;
matrice limitée à Moodle 4.1/PHP 8.0 et 4.5/PHP 8.2, sans Moodle 5.x ni PHP 8.3/8.4, alors que
\code{requires} annonce encore Moodle 4.0, non testé.

\section{Documentation et internationalisation}

\begin{itemize}
  \item \corrige{} : noms et adresses réels retirés des valeurs par défaut
    (\fichier{settings.php} supprimé, défauts vides dans \fichier{scopesettings.php}) ;
    \code{embargo\_help} dit « au plus un an ». Les anciennes valeurs restent dans l'historique
    \code{git} d'un dépôt public.
  \item README \present{} et désormais contradictoire avec le code : « 90 + 7 tests » (233),
    badge « PHP 8.2 » (8.0 et 8.2 testés), embargo de « 5 ans » (12 mois au plus,
    \code{THESE\_EMBARGO\_MAX\_MONTHS}), rapport « rédigé dans un formulaire » (déposé depuis
    0.23.0).
  \item INSTALL \partiel{} : embargo à jour, section « Rapport du directeur : rédigé sur sa propre
    page » obsolète (213--215), pas de renvoi vers CONTRIBUTING.
  \item Guide : 11 captures sur 97 non citées (2,4~Mo) ; 15~Mo d'images PNG dans le dépôt.
  \item Sept formats de date en dur (\code{\%d/\%m/\%Y}, \code{\%B \%Y}) dans
    \fichier{doctemplates.php}, \fichier{planning.php}, \fichier{planning\_page.php},
    \fichier{jurors.php}.
\end{itemize}

\section{Plan d'action proposé}

Le plan reprend celui du 25 octobre, réordonné ; les actions déjà faites sont retirées.

\begin{longtable}{@{}p{0.6cm}p{11.2cm}p{1.6cm}p{1.6cm}@{}}
\toprule
\textbf{N\textsuperscript{o}} & \textbf{Action} & \textbf{Effort} & \textbf{Impact} \\
\midrule
\endhead
1 & Activer la protection de \code{main} (fusion obligatoire, quatre contrôles requis, une
    approbation) ; contrôle de la première section de CHANGES dans le job « Publication » ;
    fichier \fichier{LICENSE}. & Petit & Élevé \\
2 & Directeur extérieur : demande émise seulement après validation du directeur interne ou de la
    DEVE, adresse affichée au validateur, exclusion des relances groupées avant validation ;
    jeton exigé dans « Mes signatures » ; tests de non-régression. & Petit & Élevé \\
3 & \code{RISK\_XSS} sur \code{managedeve} ; retrait des balises \code{<tcpdf>} ;
    substitution des balises en une passe ; passer le dossier aux derniers envois faits sans
    lui, par cohérence (le titre reste en clair pour l'université, le jury et le directeur
    extérieur, qui doivent le connaître). & Petit & Moyen \\
4 & Confidentialité : compléter métadonnées, requêtes et destinations externes ; rédiger la
    politique de conservation du PV (anonymisation plutôt que suppression). & Moyen & Élevé \\
5 & Sauvegarde : \code{templatehash}, annotation de \code{validatedby}, \code{is\_samesite} pour
    \code{jurorid} et \code{questionid}, journal des dossiers ignorés ; documenter
    \code{these\_a6} et l'annuaire ; test d'aller-retour complet (les deux modules). & Petit &
    Moyen \\
6 & Étape de mise à niveau : index manquants, clé étrangère \code{these\_question.theseid},
    adresse normalisée, \code{these\_juror.email} à 255 ; figer les valeurs de l'étape
    2026102504. & Petit & Moyen \\
7 & Suivi multi-promotions : charger étudiants et directeurs en une requête
    (\code{get\_in\_or\_equal}) dans \code{thesesynthesis\_get\_dossiers} et
    \code{thesesynthesis\_search}. & Petit & Moyen \\
8 & Horodatage d'envoi pour \code{send\_digest} ; trace par envoi dans les deux tâches. & Petit
    & Faible \\
9 & Faibles : \code{s()} dans les redirections, \code{is\_enrolled} dans l'import CSV, préfixe
    des cellules dangereuses dans les exports, jeton hors de \code{set\_url}, manuscrit limité au
    PDF, \code{maxbytes} sur les imports tableur, contrôle de \code{section}. & Petit & Moyen \\
10 & CI : concurrence conditionnelle, permissions du job guide, \code{moodle-cs} épinglé,
    \code{timeout-minutes}, ligne Moodle 5.x/PHP 8.3, plancher \code{requires} aligné sur 4.1,
    étape \code{grunt}, couverture en artefact. & Petit & Moyen \\
11 & README et INSTALL mis à jour (tests, badges, embargo, rapport déposé) ; élagage des
    captures ; formats de date localisés. & Petit & Faible \\
12 & Supprimer le code mort ; sortir les constantes de \fichier{locallib/} ; dans
    \fichier{classes/}, appeler les classes plutôt que les façades ; helper d'amorce de page et
    classe de base des formulaires. & Moyen & Moyen \\
13 & Scinder \code{signatures}, \code{doctemplates} et \code{jury} ; découper
    \code{build\_document\_context} ; sortir les contrôleurs \code{*\_page} ; créer
    \code{mod\_thesesynthesis\textbackslash local} et supprimer la dépendance inverse. & Grand &
    Moyen \\
14 & Notifications par l'API de messagerie, événements Moodle. & Grand & Moyen \\
\bottomrule
\end{longtable}

\paragraph{Recommandation de méthode.}
Les actions 1 à 3 tiennent en une demande de fusion chacune et ferment tous les constats de
sécurité moyens (le seul restant est celui du directeur extérieur, action 2). Il est conseillé de les livrer avant toute nouvelle fonctionnalité, puis de
réserver une demande de fusion de dette sur trois aux actions 4 à 13 : au rythme actuel, la
dette progresse plus vite qu'elle ne se résorbe.

\end{document}
